\section{Justificación}\label{sec:justificacion}

FIRA pidió que la estimación sirva para decisiones productivas, de planeación y eventualmente financieras, y que cada
equipo explique qué significa su error en términos productivos~\cite{preguntas}. Para cualquiera de esas decisiones un
número solo no basta. Hace falta saber cuánto puede fallar.

Las mediciones de campo son escasas y costosas, mientras que las imágenes de satélite y el clima son públicos y cubren
cualquier parcela. Vale la pena saber cuánta información del rendimiento contienen realmente. Con pocas parcelas y un
registro heterogéneo, también vale la pena medir el error con cuidado, porque un error subestimado haría parecer
seguras decisiones que no lo son.

Este trabajo aporta una estimación por parcela con un error medido de forma conservadora, su traducción a pesos y una
aplicación para consultarla. El detalle de cada decisión está en los reportes técnicos del repositorio~\cite{repo}.
